\chapter{Methodology}

\section{What timing derived numbers can and cannot claim}

The target runs inside WSL2, a Hyper-V utility virtual machine, which exposes no
hardware performance monitoring unit. Every number in this report is therefore
derived from wall clock timing and a known transfer size or operation count.
This is the classical STREAM methodology and it is entirely sufficient for
latency, bandwidth, and throughput. What it cannot provide is a direct cycle or
cache miss count from the hardware; a \texttt{perf\_event} wrapper is compiled
but is gated to bare metal Linux, where \texttt{is\_available} returns true, and
it stays dormant under WSL2.

\section{The clock does not track turbo}

Cycle denominated figures need the frequency the core actually ran at. Under
WSL2 the \texttt{cpu MHz} field in \texttt{/proc/cpuinfo} is pinned to a base
value near \SI{3.42}{\giga\hertz} and never moves, even under full load, and no
\texttt{cpufreq} interface exists. The pinned core, however, boosts to about
\SI{5.3}{\giga\hertz}. Using the base value would understate every cycle count
by roughly a factor of \num{1.6}.

The fix is to measure the running clock in software. A dependent chain of
integer additions executes at one cycle per iteration, because each addition
depends on the previous result and the loop counter runs on a separate port.
Timing a known iteration count gives the frequency directly. Because the chain
cannot run faster than one cycle per iteration, any perturbation only lengthens
the time, so the maximum frequency over several warmed passes is the cleanest
estimate. This method reports about \SI{5.3}{\giga\hertz} under load, and with
it the L1 latency lands at exactly the expected 5 cycles.

\section{Turbo stability and the power plan}

For a stable clock the Windows power plan maximum processor state is set to 99
percent, which caps turbo so the boost frequency does not wander. The software
clock measurement is taken during every run regardless, so a session without the
cap is still honest, only noisier.

\section{Pinning and the hybrid layout}

Benchmarks pin to a single logical CPU with \texttt{sched\_setaffinity}, a P
core by default. The i7-14700K is a hybrid part with 8 performance cores and 12
efficiency cores, but WSL2 does not expose the \texttt{cpu\_core} and
\texttt{cpu\_atom} split, and it reports a uniform L2 size across all cores, so P
and E cannot be recovered from sysfs. The topology reader detects the split when
present and otherwise falls back to the documented Intel enumeration, marking the
result as heuristic rather than measured.

\section{Reproducibility}

Every run records the compiler and flags, kernel version, WSL flag, measured
clock, pinned core, and git commit into the summary, so any figure can be traced
to the conditions that produced it. The reproducibility block for the run behind
this report is in Table~\ref{tab:repro}.

\begin{table}[h]
\centering
\caption{Reproducibility record for the committed run.}
\label{tab:repro}
\begin{tabular}{ll}
\hline
Compiler & g++ (Ubuntu 15.2.0-16ubuntu1) 15.2.0 \\
Kernel & 6.18.33.2-microsoft-standard-WSL2 \\
WSL2 & True \\
Measured clock & 5.34 GHz \\
Git commit & 3065706 \\
\hline
\end{tabular}

\end{table}
